% SPDX-License-Identifier: GPL-3.0-only
% Licensed under GNU GPL version 3; https://www.gnu.org/licenses/gpl-3.0.html . See LICENSE at repository root.
\section{Counterexamples, ranks, and local trace obligations}\label{sec:progress}

Ordinary rule soundness is insufficient for a cyclic proof. When a conclusion
is false, the local instance must select one false premise and a world in
which all advertised trace inequalities on that selected edge hold together.
This section defines the rank used for those inequalities.

\subsection{Execution histories}
A history is a finite list of traces. For $\pi$ and a list of additional
states, define
\begin{align*}
 H(\pi,[])&=[\pi],\\
 H(\pi,s::\eta)&=\pi::H(\pi\cat[s],\eta).
\end{align*}
Thus
\[
 H(\pi,[s_1,s_2])=[\pi,\pi\cat[s_1],\pi\cat[s_1,s_2]].
\]
An execution sequence $\execseq{\pi}{\alpha}{h}$ holds when there exist
$s$, a nonempty $\eta$, and $\zeta$ such that
\[
 s::\eta\in\interp{\alpha}_M,\qquad
 \pi\odot(s::\eta)\leadsto\zeta,\qquad h=H(\pi,\eta).
\]
Its last component is the complete accumulated result trace.
\sources{\rocq{Progress.v}{execution_history},
  \rocq{Progress.v}{execution_sequence}, \rocq{Progress.v}{execution_history_result}}

\subsection{Counterexample histories and their minimum}
The generic relation \code{overlap_list} also applies to histories: it merges
one equal endpoint \emph{trace}, just as execution overlap merges one state.
Write $h\odot h'\leadsto h''$ for this history-level use.

\begin{definition}[Counterexample history]
The relation $\ce{A}{\pi}{h}$ is defined by recursion on the outer formula
constructor. For a primitive box,
\begin{align*}
 \ce{[\alpha]A}{\pi}{h}\quad\Longleftrightarrow\quad
 \exists h_0,\zeta,h_1.\;&\execseq{\pi}{\alpha}{h_0}
   \land\mathsf{last}(h_0)=\zeta\\
   &{}\land\ce{A}{\zeta}{h_1}\land h_0\odot h_1\leadsto h.
\end{align*}
For every formula whose primitive root is not a box,
\[
 \ce{A}{\pi}{h}\quad\Longleftrightarrow\quad\pi\nsat A\land h=[\pi].
\]
\end{definition}
\sources{\rocq{Progress.v}{counterexample}}

This definition follows nested outer boxes. It does not recursively inspect
Boolean or suffix structure when collecting a history. Diamonds and dual
suffixes have primitive negation roots and therefore use the singleton-history
case. The choice of primitive representation consequently matters to the
rank definition even though the derived semantic laws remain available.

\begin{lemma}[Histories characterize failure]
For $\pi\ne[]$,
\[
 \pi\nsat A\quad\Longleftrightarrow\quad\exists h.\ \ce{A}{\pi}{h}.
\]
Every such history begins with $\pi$.
\end{lemma}
\sources{\rocq{Progress.v}{counterexample_exists},
  \rocq{Progress.v}{counterexample_false}, \rocq{Progress.v}{counterexample_head}}

\begin{definition}[Measure and rank]
The history measure is the sum of its component lengths,
\[
 \measure{[\pi_0,\ldots,\pi_n]}=\sum_{j=0}^{n}|\pi_j|.
\]
Let $C(A,\pi)=\{\measure{h}\mid\ce{A}{\pi}{h}\}$. Define
\[
 \rankf{A}{\pi}=\begin{cases}
 \min C(A,\pi),&C(A,\pi)\ne\varnothing,\\
 0,&C(A,\pi)=\varnothing.
 \end{cases}
\]
\end{definition}
\sources{\rocq{Progress.v}{sequence_measure},
  \rocq{Progress.v}{minimum_counterexample}, \rocq{Progress.v}{minimum_spec}}

This is a classical mathematical minimum in $\N$, not an executable search
through program runs. Failure on a nonempty trace gives an attained minimum
and its leastness property. The fallback zero handles an empty counterexample
set and does not replace this existence proof.

\begin{lemma}[Non-box and atomic-box ranks]
If $A$ has no primitive box root, $\pi\ne[]$, and $\pi\nsat A$, then
$\rankf{A}{\pi}=|\pi|$.
If a false atomic box has its uniquely determined extension $\pi'$---an
assignment, or a successful test---then
\[
 \rankf{[\alpha]A}{\pi}=|\pi|+\rankf{A}{\pi'}.
\]
In the test case, falsity of the box also entails success of its guard.
\end{lemma}
\sources{\rocq{Progress.v}{nonbox_minimum}, \rocq{Progress.v}{atomic_minimum}}

\begin{proof}[Proof sketch]
The non-box history is exactly $[\pi]$. For an atomic box, execution history
is $[\pi,\pi']$. Overlapping it with a body counterexample gives
$\pi::h'$. Compare an attained source minimum with an attained body minimum
in both directions to obtain equality.
\end{proof}

\subsection{Decrease and transport}
On false admissible principals the supported inequalities are:
\begin{center}
\begin{tabular}{@{}lll@{}}
\toprule
Source & Selected continuation & Rank relation\\
\midrule
Assignment box at $\pi$ & Body at $\pi\cat[s']$ & $<$\\
Successful test box at $\pi$ & Body at $\pi\cat[s]$ & $<$\\
Sequence box & Nested boxes at the same $\pi$ & $\le$\\
Choice box & One false branch at the same $\pi$ & $\le$\\
Star box & One false unfolding at the same $\pi$ & $\le$\\
Suffix or dual suffix at $s::\eta$ & False same operator at $\eta\ne[]$ & $<$\\
\bottomrule
\end{tabular}
\end{center}
\sources{\rocq{Progress.v}{assignment_decrease},
  \rocq{Progress.v}{test_decrease}, \rocq{Progress.v}{sequence_nonincrease},
  \rocq{Progress.v}{choice_nonincrease}, \rocq{Progress.v}{star_nonincrease},
  \rocq{Progress.v}{suffix_decrease}, \rocq{Progress.v}{dual_suffix_decrease}}

For assignment/test, the accumulated trace grows, while removing the outer
box removes the positive $|\pi|$ contribution from its minimum history.
For suffixes, both ranks are lengths and removing the head shortens the trace.
For choice and star, the false branch is selected from a \emph{minimum}
source counterexample. Selecting an arbitrary false branch would not justify
the required upper bound.

An occurrence rank is
\[
 \rho_w(\ell,A,\tau)=\rankf{A}{\eval{\ell}{w}}.
\]
The judgment rank $\rho_{J,w}(t)$ looks up the succedent occurrence tagged
$t$ and uses its occurrence rank; absent tags receive zero. Well-formedness
ensures uniqueness of the lookup.
\sources{\rocq{Progress.v}{occurrence_rank}, \rocq{Progress.v}{rank_in_context},
  \rocq{Rules.v}{rank_context_occ}}

\subsection{What a local rule instance proves}
An edge has a finite set $V_e$ of valid tag pairs and a subset $P_e$ of
progressing pairs. A pair $(a,b)$ connects source tag $a$ to premise tag $b$.
Respect for an edge from $J,w$ to $J',w'$ means
\begin{align*}
 &(a,b)\in V_e\ \Longrightarrow\ \rho_{J',w'}(b)\le\rho_{J,w}(a),\\
 &(a,b)\in P_e\ \Longrightarrow\ \rho_{J',w'}(b)<\rho_{J,w}(a).
\end{align*}
The source and target tags must exist at their endpoints, and $P_e\subseteq V_e$.
\sources{\rocq{../Cyclic/Framework.v}{edge_info},
  \rocq{../Cyclic/Framework.v}{respects_edge},
  \rocq{../Cyclic/Framework.v}{edge_info_wfb}}

Each rule instance stores both ordinary validity preservation and a
countermodel operation:
\[
 w\nsat J\ \Longrightarrow\
 \exists k,J_k,w'.\ J_k\text{ is premise }k\ \land\
 w'\nsat J_k\ \land\ e_k\text{ is respected from }(J,w)\text{ to }(J_k,w').
\]
All pairs on the selected edge must be respected by the \emph{same} $w'$.
Rules cannot justify each pair using a different countermodel.
\sources{\rocq{../Cyclic/Framework.v}{rule_instance},
  \rocq{Rules.v}{local_rule}, \rocq{Rules.v}{counter_rule}}

Unchanged succedent occurrences contribute identity pairs with equal rank.
The enhanced right box rules add their proved principal pairs when both
principal and continuation are tagged. Assignment/test mark those pairs
progressing; sequence, choice, and star do not. The continuing second
suffix/dual-suffix right edge can mark a strict principal pair. The first
suffix-right edge cannot. Atoms, diamonds, and untagged principals do not
gain a trace merely because a decomposition is semantically valid.
\sources{\rocq{Rules.v}{retained_pairs}, \rocq{Rules.v}{continuation_pairs},
  \rocq{Rules.v}{enhance_right}, \rocq{Rules.v}{suffix_right},
  \rocq{Rules.v}{dual_suffix_right}}

Label substitution transports ranks through its corresponding environment
change. Fresh eigenvariable switching transports histories by an involution
and preserves every component length. Tag renaming follows the same labelled
formula with its new identifier. Backlinks and semantic conversions add no
progress by themselves.
